\scanpage{3}
Pour EGA VI

\uline{Recollements d'objets d'un site}.

Soit $S$ un site\inserted{ standard (i.e. à topologie moins fine que la canonique}, muni d'une partie $\underline{M}$ de l'ensemble des flèches, satisfaisant les conditions habituelles :

a) les flèches de $M$ sont quarrab\deleted{k}\handwritten{l}es et stables par changement de base,

b) $M$ contient les flèches\handwritten{\ensuremath{\vert}}identiques et est stable par composition,

c) une flèche $u:X\to Y$ appartient à $M$ si elle y appartient localement sur $Y$,

et en plus une condition de saturation:

d) Soit $u_{\handwritten{i}}:X_{\handwritten{i}}\to\deleted{Y}\handwritten{X}$ un\handwritten{e}\inserted{\handwritten{ famille couvrante de}} \deleted{$Y$-}morphisme\handwritten{s} dans $\deleted{S}\handwritten{M}$, \deleted{supposons \handwritten{\unsure{les}} $X_{\handwritten{i}}\unsure{\to|\,} \unsure{Y|X}$ et \unsure{les|dans}}\inserted{\handwritten{(i.e, par a) et c), les}} $X_{\handwritten{i}}\times_{\deleted{Y}\handwritten{X}}X_{\handwritten{j}}\xrightarrow{\mathrm{pr}_{1}}X_{\handwritten{i}}$ dans $M$\handwritten{)}, \inserted{\handwritten{et supposons les $uu_{\unsure{i}}:X_i\to Y$ dans $M$,}} alors $u\inserted{\handwritten{{}:X\to Y}}$ est dans $M$.

% slanted
\margin{\handwritten{essayer de ne pas le mettre ici, en saturant plus tard en $\widetilde M$, et formulant $\tilde e$ de façon un peu différente : $u$ est dans $M$ sss on peut trouver des $u_i:X_i\to X$ dans $M$ avec $uu_i$ \unsure{ds|dans} $M$.}}

\uline{Proposition} I\handwritten{l} existe une partie $\overset{\handwritten{\sim}}{M}$ de l'ensemble des flèches du topos $\overset{\handwritten{\sim}}{S}$, déterminée de façon unique par les propriétés suivantes: $\overset{\handwritten{\sim}}{M}$ satisfait aux conditions \inserted{\handwritten{$\tilde a,\tilde b,\tilde c,\tilde d$}} analogues à a) à d), et de plus aux conditions suivantes:

$\overset{\handwritten{\sim}}{e}$) Si $u:X\to Y$ est une flèche de $S$, alors $u$ est dans $M$ sss $u$ est dans $\overset{\handwritten{\sim}}{M}$.

$\overset{\handwritten{\sim}}{f}$) Si $u:F\to X$ est une flèche dans $\overset{\handwritten{\sim}}{\mathbb{M}}$, avec $X\handwritten{\in}\deleted{S}\handwritten{\mathrm{Ob}\,S}$, alors il existe une famille couvrante $u_i:X_i\to F$ de flèches dans $\overset{\handwritten{\sim}}{M}$, avec $X_i\handwritten{\in}S$\handwritten{.} \deleted{et $uu_i\handwritten{\unsure{\in}}\overset{\handwritten{\sim}}{M}$.}

Démontrons l'unicité, et montrons que l'on peut décider sans ambiguité si une flèche $u:F\to G$ de $\overset{\handwritten{\sim}}{S}$ est dans $\overset{\handwritten{\sim}}{M}$. Pour ceci, distinguons deux cas:

1º) Cas où le morphisme diagonal $F\to F\times_G\deleted{\unsure{G|E}}F$ est représentable. \deleted{\unsure{U|X}\illegible{1}}

On sait qu'on peut trouver une famille couvrante $X_i\to F$, les $X_i\handwritten{\in} S$,

% archive sheet number: 2
